\section{问题重述}

新鲜药材含水量高，需经热风烘干以便贮存。烘干时把药材放入烘房，热风一方面把药材加热，一方面带走从药材表面蒸发的水分；随着水分向外迁移并蒸发，药材内部的温度和含水率不断变化，有的药材还会因失水而收缩。工程上关心两件事：药材内部温度与含水率如何随时间和位置变化，以及需要烘多久才能把水分降到规定标准。

把一根药材理想化为半径 $R_0=\SI{2}{cm}$、长 $L=\SI{25}{cm}$ 的圆柱。初始时药材温度为 $\SI{28}{\celsius}$、干基含水率为 $2.55$，且内部均匀。烘房的空气温度 $\Tair(t)$ 与环境含水浓度 $\Cenv(t)$ 随时间给定（附件 1，$0$ 至 $\SI{4}{h}$，每 $\SI{60}{s}$ 一个数据点）。药材表面与热风之间的对流换热系数为 $h=\SI{25}{W/(m^2.K)}$，传质系数为 $h_m=\SI{8e-7}{m/s}$。在此基础上，题目要求解决以下四个问题。

\begin{description}[leftmargin=2.4em,itemsep=2pt]
\item[问题一] 在附录 2 给定的一组物性参数下，求预热阶段（$0$ 至 $\SI{1800}{s}$）药材内部的温度场与含水率场，并按每 $\SI{1}{s}$、每 $\SI{0.1}{cm}$ 输出（表~1、表~2 与 result1）。
\item[问题二] 改用附录 3 给定的、随状态变化的物性参数，求整个干燥过程的温度场与含水率场（表~3、表~4 与 result2）。
\item[问题三] 以「药材各处含水率均低于 $0.15$」为达标标准，求达标所需的干燥时长（单位 h），并输出达标前的含水率场（表~5 与 result3）。
\item[问题四] 考虑药材在干燥中因失水而收缩，半径随时间的变化由附件 2 给定；在附录 4 的物性参数下重新求解并给出达标时长（表~6 与 result4）。
\end{description}

四个问题层层递进：问题一是常物性下的预热特例；问题二把物性换成随状态变化并延伸到全过程；问题三在问题二的解上确定达标时刻；问题四进一步引入几何收缩。所有结果均保留四位小数。
